\documentclass{article}
\usepackage{amsfonts} 
\title{Longitud de Arco de una Curva Regular en el Espacio Tridimensional}
\author{Por: David Duarte}
\date{\today}

\begin{document}
	
	\maketitle
	
	\section{Longitud de arco de una curva regular en el espacio tridimensional}

\subsection {Introducción:}

En la geometría diferencial, la longitud de arco de una curva en el espacio tridimensional es una medida fundamental que nos permite cuantificar la "longitud" de una curva en un espacio tridimensional.
La longitud de arco de una curva paramétrica en ($\mathbb{R}^3$) es una generalización del concepto de longitud de arco en una curva en el plano.

\subsection{Definición de curva regular en $\mathbb{R}^3:$}

Una curva paramétrica ($\mathbf{r}(t) = (x(t), y(t), z(t))$) definida en un intervalo ($a \leq t \leq b$) se dice que es una curva regular si la derivada ($\mathbf{r}'(t)$) es continua y nunca es el vector cero para todo (t) en el intervalo.
	
	\section{Definición de Curva Regular en $\mathbb{R}^3$}
	Una curva paramétrica $\mathbf{r}(t) = (x(t), y(t), z(t))$ definida en un intervalo $a \leq t \leq b$ se dice que es una curva regular si la derivada $\mathbf{r}'(t)$ es continua y nunca es el vector cero para todo $t$ en el intervalo.
	
	\section{Cálculo de la Longitud de Arco}
	La longitud de arco de una curva paramétrica en $\mathbb{R}^3$ se puede calcular utilizando una integral definida. La fórmula para la longitud de arco de una curva paramétrica $\mathbf{r}(t)$ en el intervalo $[a, b]$ es dada por:
	\[ L = \int_{a}^{b} \| \mathbf{r}'(t) \| \, dt \]
	donde $\| \mathbf{r}'(t) \|$ representa la magnitud del vector derivada de la curva.
	
	\section{Aplicaciones de la Longitud de Arco}
	La longitud de arco de una curva paramétrica es útil en diversas aplicaciones, como en física, ingeniería, geometría computacional y modelado 3D.
	

	
\end{document}

